% Experiment A LaTeX section
% Requires: \usepackage{graphicx}
% The default figure root assumes this file is compiled from its
% containing analysis directory. Override \ExperimentAFigureRoot
% in the main report if the figures are copied elsewhere.
\providecommand{\ExperimentAFigureRoot}{representative_examples}

\section{Experiment A: Zero-Shot Grounding DINO on Thermal Imagery}
\label{sec:experiment-a}

\subsection{Research Question}

This experiment evaluates how well pretrained zero-shot Grounding DINO transfers to BirdsAI thermal imagery without training or fine-tuning.

\subsection{Experimental Setup}

The experiment used \texttt{IDEA-Research/grounding-dino-tiny} through the Hugging Face Transformers implementation. Custom CUDA and CPU kernels were disabled. Predictions were mapped to the two BirdsAI categories, \emph{animal} and \emph{person}. Ambiguous phrases matching both categories were excluded.

Four prompt formulations were compared on the validation split. The prompt was selected using COCO mAP@[0.50:0.95], and the operating score threshold was selected using F1 at IoU 0.50. The frozen configuration used the prompt \texttt{person . animal .}, a box threshold of 0.10, a text threshold of 0.20, and an operating score threshold of 0.10. TestReal was evaluated once after configuration selection.

\subsection{Dataset Distribution}

\begin{table}[htbp]
\centering
\caption{Validation and TestReal distribution comparison.}
\label{tab:experiment-a-distribution}
\small
\begin{tabular}{lrrrrrrr}
\hline
Split & Seq. & Images & Ann. & Ann./image & Empty & Animal area & Person area \\
\hline
Validation & 6 & 6735 & 15086 & 2.24 & 54.2\% & 294.0 & 300.0 \\
Test & 16 & 21333 & 78761 & 3.69 & 27.4\% & 2952.0 & 522.0 \\
\hline
\end{tabular}
\end{table}

TestReal contains substantially larger animals, more annotations per image, and fewer empty images than validation. In particular, median animal bounding-box area increases from 294 pixels squared in validation to 2952 pixels squared in TestReal. This distribution difference makes TestReal more favorable to zero-shot object detection and helps explain the higher test performance.

\subsection{Validation Prompt Comparison}

\begin{table}[htbp]
\centering
\caption{Corrected validation results after excluding ambiguous phrases.}
\label{tab:experiment-a-prompts}
\small
\begin{tabular}{lrrrrrr}
\hline
Prompt & mAP & AP50 & AP75 & F1@0.25 & Animal AP50 & Person AP50 \\
\hline
\texttt{person\_animal} & 1.16\% & 3.36\% & 0.49\% & 12.29\% & 6.72\% & 0.00\% \\
\texttt{human\_animal} & 0.38\% & 0.78\% & 0.32\% & 3.17\% & 1.56\% & 0.00\% \\
\texttt{person\_wildlife} & 0.75\% & 2.26\% & 0.40\% & 12.72\% & 4.53\% & 0.00\% \\
\texttt{person\_species} & 0.66\% & 1.17\% & 0.64\% & 2.38\% & 2.34\% & 0.00\% \\
\hline
\end{tabular}
\end{table}

The \texttt{person\_animal} prompt achieved the highest validation COCO mAP and was therefore selected before TestReal was examined.

\subsection{Final Results}

\begin{table}[htbp]
\centering
\caption{Final validation and TestReal results for the frozen configuration.}
\label{tab:experiment-a-final}
\small
\begin{tabular}{lrrrrrrrr}
\hline
Split & mAP & AP50 & AP75 & Precision & Recall & F1 & FP/frame & Person AP50 \\
\hline
Validation & 1.16\% & 3.36\% & 0.49\% & 15.71\% & 12.85\% & 14.13\% & 1.54 & 0.00\% \\
Test & 5.49\% & 18.24\% & 1.97\% & 33.40\% & 35.22\% & 34.29\% & 2.59 & 6.00\% \\
\hline
\end{tabular}
\end{table}

\begin{table}[htbp]
\centering
\caption{TestReal class-specific results at score threshold 0.10.}
\label{tab:experiment-a-classes}
\small
\begin{tabular}{lrrrrrrr}
\hline
Category & TP & FP & FN & Precision & Recall & F1 & AP50 \\
\hline
Animal & 23917 & 34951 & 32991 & 40.63\% & 42.03\% & 41.32\% & 30.49\% \\
Person & 3824 & 20357 & 18029 & 15.81\% & 17.50\% & 16.61\% & 6.00\% \\
\hline
\end{tabular}
\end{table}

On TestReal, the model achieved an overall mAP of 5.49\%, AP50 of 18.24\%, and F1 of 34.29\%. Animal detection transferred substantially better than person detection: animal AP50 was 30.49\%, whereas person AP50 was 6.00\%. The low AP75 of 1.97\% indicates that bounding-box localization remained poor. The model also produced 2.59 false positives per frame.

\subsection{Qualitative Examples}

Figures~\ref{fig:experiment-a-animal-examples} and \ref{fig:experiment-a-person-examples} show representative true positives, false positives, and false negatives. Green boxes indicate ground truth and red boxes indicate predictions.

\begin{figure}[htbp]
\centering
\begin{minipage}[t]{0.31\textwidth}
\centering
\includegraphics[width=\linewidth]{\ExperimentAFigureRoot/tp/animal/tp_animal_01_image_7431.jpg}

\small TP Animal
\end{minipage}
\hfill
\begin{minipage}[t]{0.31\textwidth}
\centering
\includegraphics[width=\linewidth]{\ExperimentAFigureRoot/fp/animal/fp_animal_01_image_10758.jpg}

\small FP Animal
\end{minipage}
\hfill
\begin{minipage}[t]{0.31\textwidth}
\centering
\includegraphics[width=\linewidth]{\ExperimentAFigureRoot/fn/animal/fn_animal_01_image_6549.jpg}

\small FN Animal
\end{minipage}

\caption{Representative animal detection outcomes on TestReal.}
\label{fig:experiment-a-animal-examples}
\end{figure}

\begin{figure}[htbp]
\centering
\begin{minipage}[t]{0.31\textwidth}
\centering
\includegraphics[width=\linewidth]{\ExperimentAFigureRoot/tp/person/tp_person_01_image_12565.jpg}

\small TP Person
\end{minipage}
\hfill
\begin{minipage}[t]{0.31\textwidth}
\centering
\includegraphics[width=\linewidth]{\ExperimentAFigureRoot/fp/person/fp_person_01_image_6434.jpg}

\small FP Person
\end{minipage}
\hfill
\begin{minipage}[t]{0.31\textwidth}
\centering
\includegraphics[width=\linewidth]{\ExperimentAFigureRoot/fn/person/fn_person_01_image_12941.jpg}

\small FN Person
\end{minipage}

\caption{Representative person detection outcomes on TestReal.}
\label{fig:experiment-a-person-examples}
\end{figure}

\subsection{Limitations}

\begin{itemize}
\item Only one pretrained Grounding DINO model size was evaluated.
\item No thermal-domain fine-tuning or adaptation was performed.
\item Validation contains only six sequences and differs substantially from TestReal in object size, density, and empty-frame frequency.
\item Detection performance is strongly influenced by object size.
\item The results do not include sequence-level confidence intervals.
\end{itemize}

\subsection{Conclusion}

Zero-shot Grounding DINO demonstrates partial but limited transfer to BirdsAI thermal imagery. It detects animals substantially better than people, while precise localization remains weak. The model provides a useful zero-shot baseline and may assist pseudo-label generation, but it is not sufficiently reliable for deployment without thermal-domain adaptation or fine-tuning.
